\chapter{碱金属原子和电子自旋}
\problem{4.1}{已知 $\mathrm{Li}$ 原子光谱主线系最长波长 $\lambda=6707\,\angstrom$，辅线系系限波长 $\lambda_\infty=3519\,\angstrom$。求 $\mathrm{Li}$ 原子第一激发电势和电离电势。}
\answerarea
\problem{4.2}{$\mathrm{Na}$ 原子的基态为 $3\mathrm S$。已知其共振线波长为 $5893\,\angstrom$，漫线系第一条的波长为 $8193\,\angstrom$，基线系第一条的波长为 $18459\,\angstrom$，主线系的系限波长为 $2413\,\angstrom$。试求 $3\mathrm S,3\mathrm P,3\mathrm D,4\mathrm F$ 各谱项的项值。}
\answerarea\clearpage
\problem{4.3}{$\mathrm K$ 原子共振线波长为 $7663\,\angstrom$，主线系系限波长为 $2858\,\angstrom$。已知 $\mathrm K$ 原子基态为 $4\mathrm S$，试求 $4\mathrm S,4\mathrm P$ 谱项的量子数修正项 $\Delta_{\mathrm S},\Delta_{\mathrm P}$ 值各为多少？}
\answerarea
\problem{4.4}{$\mathrm{Li}$ 原子的基态项为 $2\mathrm S$。当把 $\mathrm{Li}$ 原子激发到 $3\mathrm P$ 态后，问当 $\mathrm{Li}$ 从 $3\mathrm P$ 激发态向低能级跃迁时可能产生哪些谱线（不考虑精细结构）？}
\answerarea\clearpage
\problem{4.5}{为什么谱项 $\mathrm S$ 项的精细结构总是单层结构？试直接从碱金属光谱双线的规律性和从电子自旋与轨道相互作用的物理概念两方面分别说明之。}
\answerarea
\problem{4.6}{试计算氢原子赖曼线系第一条的精细结构分裂的波长差。}
\answerarea\clearpage
\problem{4.7}{$\mathrm{Na}$ 原子光谱中得知其 $3\mathrm D$ 项的项值 $T_{3\mathrm D}=1.2274\times10^6\,\mathrm{m^{-1}}$，试计算该谱项之精细结构裂距。}
\answerarea
\problem{4.8}{原子在热平衡条件下处在各种不同能量激发态的原子的数目是按玻耳兹曼分布的，即能量为 $E$ 的激发态原子数目
\[
N=N_0\frac{g}{g_0}\exp\left(-\frac{E-E_0}{kT}\right),
\]
其中 $N_0$ 是能量为 $E_0$ 的状态的原子数，$g$ 和 $g_0$ 是相应能量状态的统计权重，$k$ 是玻耳兹曼常数。从高温铯原子气体光谱中测出其共振光谱双线 $\lambda_1=8943.5\,\angstrom$、$\lambda_2=8521.1\,\angstrom$ 的强度比 $I_1:I_2=2:3$。试估算此气体的温度。已知相应能级的统计权重 $g_1=2,g_2=4$。}
\answerarea\clearpage
